\documentclass[11pt]{article}
% Standalone preamble for web-only supplementary notes.
\usepackage[T1]{fontenc}
\usepackage{amsmath,amssymb,amsthm}
\usepackage{hyperref}
\usepackage{cleveref}
\DeclareMathOperator{\Map}{Map}
\DeclareMathOperator{\Fun}{Fun}
\newcommand{\iso}{\xrightarrow{\sim}}
\newtheorem{lemma}{Supplementary Lemma}
\newtheorem{corollary}[lemma]{Supplementary Corollary}
\crefname{lemma}{Supplementary Lemma}{Supplementary Lemmas}
\Crefname{lemma}{Supplementary Lemma}{Supplementary Lemmas}
\crefname{corollary}{Supplementary Corollary}{Supplementary Corollaries}
\Crefname{corollary}{Supplementary Corollary}{Supplementary Corollaries}

\begin{document}
\section*{Constant diagrams and the exponential law}

The proof that \href{groupoids.html}{functor categories with groupoid target
are groupoids} uses compatibility of the exponential law with constant
diagrams. We describe the comparison of the actual functors and the
product calculation giving it.

For absolute categories $T,X$, let
$\Delta_T^X\colon X\to\Fun(T,X)$ be the curry of
$\pi_1\colon X\times T\to X$. Write
\[
 E_{T,C,X}\colon\Fun(T,\Fun(C,X))\longrightarrow\Fun(T\times C,X)
\]
for the chosen exponential equivalence. The first three statements use
the Chapter 1 mapping-anima and functor-category structure.

\begin{lemma}
\label[lemma]{lem:Supplementary_Constant_Naturality}
For $F\colon C\to D$, one has an identification of functors
\[
 \Fun(T,F)\circ\Delta_T^C=\Delta_T^D\circ F.
\]
\end{lemma}
\begin{proof}[Proof idea]
Uncurry both sides. For the left side, postcomposition compatibility and
the currying computation comparison identify the result with $F\pi_1$.
For the right side, restriction compatibility gives
$\pi_1(F\times\mathrm{id}_T)$, which the product projection comparison
also identifies with $F\pi_1$. Reflect through uncurrying.
\end{proof}

\begin{lemma}
\label[lemma]{lem:Supplementary_Constant_Exponential}
Put $N=\Fun(C,X)$ and let $\sigma\colon C\times T\to T\times C$
exchange the factors. The following functors $N\to\Fun(C\times T,X)$
are identified:
\[
 \sigma^*\circ E_{T,C,X}\circ\Delta_T^N
   =E_{C,T,X}\circ\Fun(C,\Delta_T^X).
\]
\end{lemma}
\begin{proof}[Proof idea]
Uncurry once. Both functors become evaluation $N\times C\to X$
precomposed with a functor $N\times(C\times T)\to N\times C$.
For the left side, exchange $C,T$, reassociate to $(N\times T)\times C$,
and forget $T$. For the right side, reassociate to $(N\times C)\times T$
and forget $T$.

Compute their projections using the product comparisons. In both cases,
the first projection is $\pi_1$ and the second is $\pi_1\pi_2$ on
$N\times(C\times T)$. The product universal property compares the two
argument functors. Postcompose with evaluation and reflect through uncurrying.
\end{proof}

\begin{corollary}
\label[corollary]{cor:Supplementary_Constant_Equivalence}
If $\Delta_T^X$ is an equivalence, then so is $\Delta_T^{\Fun(C,X)}$.
\end{corollary}
\begin{proof}[Proof idea]
Postcomposition $\Fun(C,\Delta_T^X)$ is an equivalence, as are both
exponential functors and precomposition by $\sigma$.
Use \Cref{lem:Supplementary_Constant_Exponential} and cancel these
equivalences. For $T=[1]$, this is the comparison used in the closure
of groupoids under functor categories.
\end{proof}

\section*{The coproduct comparison}

The following calculation is used in
\href{recognizing-animae.html}{Section 2.5}. We work in that section's
setup, including coproducts. Its linked Agda module also carries the Rezk
and recognition assumptions needed for the subsequent closure results.
The calculation itself uses naturality and the coproduct universal property.

\begin{lemma}
\label[lemma]{lem:Supplementary_Constant_Coproduct}
Put $d_C=\Delta_{[1]}^C\colon C\to\operatorname{Ar}(C)$, and let
$i_C,i_D$ be the inclusions into $C\amalg D$.
For the comparison
\[
 q=[\operatorname{Ar}(i_C),\operatorname{Ar}(i_D)]\colon
 \operatorname{Ar}(C)\amalg\operatorname{Ar}(D)\to\operatorname{Ar}(C\amalg D),
\]
there is an identification
\[
 q\circ(d_C\amalg d_D)=d_{C\amalg D}.
\]
\end{lemma}
\begin{proof}[Proof idea]
Restrict the left side to each summand. The copairing computation and
precomposition comparisons give $\operatorname{Ar}(i_C)d_C$ and
$\operatorname{Ar}(i_D)d_D$. By \Cref{lem:Supplementary_Constant_Naturality},
these identify with $d_{C\amalg D}i_C$ and $d_{C\amalg D}i_D$.
The coproduct universal property assembles the two comparisons into the
displayed identification. This computes the specific factorization
used in the closure argument.
\end{proof}
\end{document}
